% =========================================================
%  ppt/09-integrasi-fsm-astar.tex
%  Materi 9 — Integrasi FSM dan A* (Mini Proyek AI)
%  Bab XI buku ajar
% =========================================================
\documentclass[aspectratio=169,10pt]{beamer}
% =========================================================
%  ppt/config/packages.tex
%  Paket bersama untuk seluruh presentasi kuliah
%  Pengantar Pemrograman Game
% =========================================================

\usepackage[T1]{fontenc}
\usepackage[utf8]{inputenc}
\usepackage[indonesian]{babel}
\usepackage{lmodern}
\usepackage{textcomp}
\usepackage{amsmath}
\usepackage{amssymb}
\usepackage{graphicx}
\usepackage{booktabs}
\usepackage{tabularx}
\usepackage{array}
\usepackage{multicol}
\usepackage{listings}
\usepackage{upquote}
\usepackage{adjustbox}
\usepackage{tcolorbox}
\usepackage{tikz}
\usetikzlibrary{
  arrows.meta,
  positioning,
  shapes.geometric,
  shapes.misc,
  calc,
  fit,
  backgrounds,
  matrix,
  chains,
  decorations.pathreplacing,
  patterns
}

\tcbuselibrary{skins}

% Jangan memuat geometry: ukuran halaman diatur oleh beamer.

% =========================================================
%  ppt/config/theme.tex
%  Satu tema Beamer untuk keempat belas presentasi
%  Pengantar Pemrograman Game
% =========================================================

% --------------------------------------------------------- warna
\definecolor{ubNavy}{HTML}{1F3864}
\definecolor{ubAccent}{HTML}{2E75B6}
\definecolor{ubHead}{HTML}{D9E2F3}
\definecolor{ubSoft}{HTML}{F2F6FC}
\definecolor{ubLine}{HTML}{BFBFBF}
\definecolor{ubText}{HTML}{1F2937}
\definecolor{ubMuted}{HTML}{5A6678}
\definecolor{ubGreen}{HTML}{1E7B45}
\definecolor{ubGreenBg}{HTML}{F1F8F2}
\definecolor{ubRed}{HTML}{B3261E}
\definecolor{ubRedBg}{HTML}{FDF2F1}
\definecolor{ubWarn}{HTML}{A35A05}
\definecolor{ubWarnBg}{HTML}{FFF8EC}
\definecolor{ubGray}{HTML}{F6F6F6}
\definecolor{ubCodeBg}{HTML}{10233F}
\definecolor{ubCodeText}{HTML}{E8EEF7}
\definecolor{ubCodeKey}{HTML}{8FBCF0}
\definecolor{ubCodeCom}{HTML}{8FB37A}
\definecolor{ubCodeStr}{HTML}{D9A46B}
\definecolor{ubCodeNum}{HTML}{7E99BE}

% Nama warna yang dipakai gambar TikZ buku ajar.
\definecolor{ubPrimary}{HTML}{1F3864}
\definecolor{ubAlt}{HTML}{F2F6FC}

% --------------------------------------------------------- dasar
\usetheme{default}
\usefonttheme{professionalfonts}
\setbeamertemplate{navigation symbols}{}
\setbeamerfont{frametitle}{size=\Large,series=\bfseries}
\setbeamerfont{framesubtitle}{size=\small}
\setbeamerfont{footline}{size=\tiny}
\setbeamerfont{block title}{size=\small,series=\bfseries}
\setbeamerfont{block body}{size=\small}
\setbeamerfont{itemize/enumerate body}{size=\small}
\setbeamerfont{itemize/enumerate subbody}{size=\footnotesize}

\setbeamersize{text margin left=0.60cm,text margin right=0.60cm}
\uchyph=0

\setbeamercolor{normal text}{fg=ubText,bg=white}
\setbeamercolor{structure}{fg=ubAccent}
\setbeamercolor{frametitle}{fg=ubNavy,bg=white}
\setbeamercolor{title}{fg=ubNavy}
\setbeamercolor{footline}{fg=ubMuted,bg=white}
\setbeamercolor{itemize item}{fg=ubAccent}
\setbeamercolor{itemize subitem}{fg=ubAccent}
\setbeamercolor{block title}{fg=white,bg=ubNavy}
\setbeamercolor{block body}{fg=ubText,bg=ubSoft}
\setbeamercolor{block title alerted}{fg=white,bg=ubWarn}
\setbeamercolor{block body alerted}{fg=ubText,bg=ubWarnBg}
\setbeamercolor{block title example}{fg=white,bg=ubGreen}
\setbeamercolor{block body example}{fg=ubText,bg=ubGreenBg}
\setbeamercolor{caption name}{fg=ubMuted}

\setbeamertemplate{blocks}[rounded][shadow=false]
\setbeamertemplate{itemize item}{\raisebox{0.10em}{\color{ubAccent}\scriptsize$\blacksquare$}}
\setbeamertemplate{itemize subitem}{\raisebox{0.06em}{\color{ubAccent}\tiny$\blacktriangleright$}}
\setbeamertemplate{enumerate item}{%
  \makebox[1.3em][l]{\color{ubAccent}\insertenumlabel.}}

\addtobeamertemplate{itemize/enumerate body begin}{%
  \setlength{\itemsep}{0.42em}\setlength{\parskip}{0pt}}{}

% --------------------------------------------------------- judul frame
\setbeamertemplate{frametitle}[default][left]

% --------------------------------------------------------- kaki
\newcommand{\materino}{00}
\newcommand{\materifoot}{Pengantar Pemrograman Game \textperiodcentered\ Materi \materino}

\setbeamertemplate{footline}{%
  \begin{beamercolorbox}[wd=\paperwidth,ht=0.015cm,dp=0cm]{footline}%
    \color{ubHead}\rule{\paperwidth}{0.4pt}%
  \end{beamercolorbox}%
  \begin{beamercolorbox}[wd=\paperwidth,ht=0.40cm,dp=0.14cm,%
                         leftskip=0.60cm,rightskip=0.60cm]{footline}%
    \usebeamerfont{footline}\color{ubMuted}\materifoot\hfill\insertframenumber%
  \end{beamercolorbox}%
}

% --------------------------------------------------------- makro buku ajar
\providecommand{\code}[1]{\texttt{\detokenize{#1}}}
\providecommand{\istilah}[1]{\emph{#1}}
\providecommand{\CPMK}[1]{CPMK\,#1}
\providecommand{\CPL}[1]{CPL\,#1}
\providecommand{\SUBCPMK}[1]{Sub-CPMK\,#1}

% --------------------------------------------------------- gaya kode
\lstdefinelanguage{CSharp}{
  morekeywords={abstract,as,base,bool,break,byte,case,catch,char,class,const,
    continue,decimal,default,delegate,do,double,else,enum,event,explicit,
    extern,false,finally,fixed,float,for,foreach,goto,if,implicit,in,int,
    interface,internal,is,lock,long,namespace,new,null,object,operator,out,
    override,params,private,protected,public,readonly,ref,return,sbyte,
    sealed,short,sizeof,stackalloc,static,string,struct,switch,this,throw,
    true,try,typeof,uint,ulong,unchecked,unsafe,ushort,using,virtual,void,
    volatile,while,var,get,set},
  sensitive=true,
  morecomment=[l]{//},
  morecomment=[s]{/*}{*/},
  morestring=[b]",
}

\lstdefinestyle{kode}{
  basicstyle=\ttfamily\footnotesize\color{ubCodeText},
  keywordstyle=\color{ubCodeKey}\bfseries,
  commentstyle=\color{ubCodeCom}\itshape,
  stringstyle=\color{ubCodeStr},
  numberstyle=\tiny\color{ubCodeNum},
  numbers=left,
  numbersep=7pt,
  backgroundcolor=\color{ubCodeBg},
  frame=none,
  showstringspaces=false,
  breaklines=false,
  tabsize=4,
  columns=fullflexible,
  keepspaces=true,
  xleftmargin=12pt,
  xrightmargin=2pt,
  aboveskip=0pt,
  belowskip=0pt,
  upquote=true,
}
\lstset{style=kode}

\newcommand{\kodejelas}[1]{%
  \vspace{0.16cm}%
  {\footnotesize\color{ubMuted}#1\par}%
}
\newcommand{\sumber}[1]{%
  \vspace{0.08cm}%
  {\scriptsize\color{ubMuted}\itshape #1\par}%
}
\newcommand{\awalan}[1]{%
  {\small\color{ubMuted}#1\par}\vspace{0.18cm}%
}

% --------------------------------------------------------- kotak
\newtcolorbox{kotak}[1][]{enhanced, colback=ubSoft, colframe=ubAccent,
  boxrule=0pt, leftrule=2.2pt, arc=1pt, left=6pt, right=6pt, top=4pt,
  bottom=4pt, fontupper=\footnotesize, #1}
\newtcolorbox{peringatan}[1][]{enhanced, colback=ubWarnBg, colframe=ubWarn,
  boxrule=0pt, leftrule=2.2pt, arc=1pt, left=6pt, right=6pt, top=4pt,
  bottom=4pt, fontupper=\footnotesize, #1}
\newtcolorbox{kiat}[1][]{enhanced, colback=ubGreenBg, colframe=ubGreen,
  boxrule=0pt, leftrule=2.2pt, arc=1pt, left=6pt, right=6pt, top=4pt,
  bottom=4pt, fontupper=\footnotesize, #1}

% --------------------------------------------------------- sampul
% \deckcover{nomor}{minggu}{judul}{sub-cpmk}{teknologi}{sumber}{blok}
\newcommand{\deckcover}[7]{%
\begin{frame}[plain]
\begin{tikzpicture}[remember picture, overlay]
  \fill[ubNavy] (current page.south west) rectangle ([xshift=6.00cm]current page.north west);
  \fill[ubAccent] ([xshift=6.00cm]current page.south west) rectangle ([xshift=6.08cm]current page.north west);

  \node[anchor=north west, text width=4.72cm, align=left, inner sep=0pt]
    at ($(current page.north west)+(0.66cm,-0.60cm)$) {%
      {\color{white}\bfseries\fontsize{10.2}{12.2}\selectfont
       Universitas Bale Bandung}\\[4pt]
      {\color{ubHead}\fontsize{7.6}{9.4}\selectfont
       Fakultas Teknologi Informasi}\\[1pt]
      {\color{ubHead}\fontsize{7.6}{9.4}\selectfont
       Program Studi Teknik Informatika}\\[7pt]
      {\color{ubAccent}\rule{4.0cm}{0.5pt}}%
    };

  \node[anchor=north west, text width=4.72cm, align=left, inner sep=0pt]
    at ($(current page.north west)+(0.66cm,-2.72cm)$) {%
      {\color{ubAccent}\fontsize{6.4}{7.8}\selectfont MATA KULIAH}\\[3pt]
      {\color{white}\bfseries\fontsize{9.8}{12.0}\selectfont
       Pengantar Pemrograman Game}\\[3pt]
      {\color{ubHead}\fontsize{7.6}{9.4}\selectfont
       MWCA009 \textperiodcentered\ 3 SKS}\\[8pt]
      {\color{ubHead}\fontsize{7.6}{9.4}\selectfont
       Semester Ganjil 2026/2027}%
    };

  \node[anchor=south west, text width=4.72cm, align=left, inner sep=0pt]
    at ($(current page.south west)+(0.66cm,0.52cm)$) {%
      {\color{ubAccent}\fontsize{6.4}{7.8}\selectfont DOSEN PENGAMPU}\\[3pt]
      {\color{white}\fontsize{8.4}{10.4}\selectfont
       Cecep Suwanda, S.Si., M.Kom.}%
    };

  \node[anchor=north west, text width=8.95cm, align=left, inner sep=0pt]
    at ($(current page.north east)+(-9.40cm,-0.92cm)$) {%
      {\color{ubAccent}\bfseries\fontsize{8.2}{10.2}\selectfont
       MATERI #1\ \ \textperiodcentered\ \ MINGGU #2}\\[8pt]
      {\color{ubNavy}\bfseries\fontsize{19.4}{23.4}\selectfont
       \hyphenpenalty=10000 \exhyphenpenalty=10000 \raggedright #3}%
    };

  \node[anchor=north west, text width=8.95cm, align=left, inner sep=0pt]
    at ($(current page.north east)+(-9.40cm,-4.30cm)$) {%
      \renewcommand{\arraystretch}{1.30}%
      \begin{tabular}{@{}>{\color{ubMuted}\fontsize{7.3}{9.0}\selectfont}p{1.85cm}
                        >{\color{ubText}\fontsize{8.1}{10.0}\selectfont}p{6.7cm}@{}}
        Sub-CPMK      & #4 \\
        Teknologi     & #5 \\
        Sumber materi & #6 \\
        Blok          & #7 \\
      \end{tabular}%
    };
\end{tikzpicture}
\end{frame}%
}

% --------------------------------------------------------- pembatas bagian
\newcommand{\sectionframe}[3]{%
\begin{frame}[plain]
\begin{tikzpicture}[remember picture, overlay]
  \fill[ubNavy] (current page.south west) rectangle (current page.north east);
  \fill[ubAccent] (current page.north west) rectangle ([xshift=0.10cm]current page.south west);
  \node[anchor=west, inner sep=0pt]
    at ($(current page.west)+(1.00cm,1.45cm)$)
    {\color{ubNavy!52!white}\bfseries\fontsize{58}{58}\selectfont #1};
  \node[anchor=west, inner sep=0pt, text width=11.6cm, align=left]
    at ($(current page.west)+(1.05cm,-0.55cm)$) {%
      {\color{ubAccent}\bfseries\fontsize{8.0}{10.0}\selectfont BAGIAN}\\[3pt]
      {\color{white}\bfseries\fontsize{20.5}{24.5}\selectfont #2}\\[5pt]
      {\color{ubHead}\small #3}%
    };
  \node[anchor=south east, inner sep=0pt]
    at ($(current page.south east)+(-0.60cm,0.32cm)$)
    {\color{ubHead}\fontsize{7.4}{9}\selectfont #1 / 6
     \ \ \textperiodcentered\ \ \insertframenumber};
\end{tikzpicture}
\end{frame}%
}

% --------------------------------------------------------- peta pertemuan
\newcommand{\blokpeta}[6]{%
  \ifnum#2=#3
    \node[rounded corners=3pt, fill=ubNavy, draw=ubNavy, line width=0.5pt,
          text width=3.02cm, minimum height=1.62cm, align=left, inner sep=5pt]
      at (#1,0) {%
        {\color{ubAccent}\bfseries\fontsize{6.8}{8}\selectfont #4}\\[1pt]
        {\color{white}\bfseries\fontsize{7.6}{9}\selectfont #5}\\[2pt]
        {\color{ubHead}\fontsize{6.4}{7.6}\selectfont #6}};
  \else
    \node[rounded corners=3pt, fill=ubSoft, draw=ubHead, line width=0.5pt,
          text width=3.02cm, minimum height=1.62cm, align=left, inner sep=5pt]
      at (#1,0) {%
        {\color{ubAccent}\bfseries\fontsize{6.8}{8}\selectfont #4}\\[1pt]
        {\color{ubNavy}\bfseries\fontsize{7.6}{9}\selectfont #5}\\[2pt]
        {\color{ubMuted}\fontsize{6.4}{7.6}\selectfont #6}};
  \fi
}

% \petapertemuan{blok aktif}{sebelumnya}{sesudahnya}{keluaran}
% \petapertemuan{blok aktif}{sebelumnya}{pertemuan ini}{sesudahnya}{keluaran}
\newcommand{\petapertemuan}[5]{%
\begin{frame}{Peta Pertemuan}
{\scriptsize\color{ubMuted}Kedudukan materi dalam empat blok perkuliahan}\par
\vspace{0.10cm}
\centering
\begin{tikzpicture}
  \blokpeta{-5.58}{1}{#1}{BLOK I}{Materi 1--5}{Game loop, vektor,\\fisika, collision}
  \blokpeta{-1.86}{2}{#1}{BLOK II}{Materi 6--9}{AI NPC: chase, FSM,\\pathfinding A*}
  \blokpeta{1.86}{3}{#1}{BLOK III}{Materi 10--11}{Unity: scene,\\level, physics}
  \blokpeta{5.58}{4}{#1}{BLOK IV}{Materi 12--14}{Aset visual, audio,\\UI, playtest}
\end{tikzpicture}
\par\vspace{0.30cm}
\begin{columns}[T,onlytextwidth]
  \begin{column}{0.315\textwidth}
    \begin{beamercolorbox}[wd=\linewidth,sep=4pt,rounded=true]{block body}
      {\color{ubMuted}\fontsize{6.6}{8}\selectfont SEBELUMNYA}\\[2pt]
      {\footnotesize #2}
    \end{beamercolorbox}
  \end{column}
  \begin{column}{0.315\textwidth}
    \begin{beamercolorbox}[wd=\linewidth,sep=4pt,rounded=true]{block title}
      {\color{ubHead}\fontsize{6.6}{8}\selectfont PERTEMUAN INI}\\[2pt]
      {\footnotesize #3}
    \end{beamercolorbox}
  \end{column}
  \begin{column}{0.315\textwidth}
    \begin{beamercolorbox}[wd=\linewidth,sep=4pt,rounded=true]{block body}
      {\color{ubMuted}\fontsize{6.6}{8}\selectfont SESUDAHNYA}\\[2pt]
      {\footnotesize #4}
    \end{beamercolorbox}
  \end{column}
\end{columns}
\par\vspace{0.26cm}
{\footnotesize\color{ubText}\textbf{Keluaran pertemuan.} #5\par}
\end{frame}%
}

% --------------------------------------------------------- penutup
\newcommand{\penutup}[2]{%
\begin{frame}[plain]
\begin{tikzpicture}[remember picture, overlay]
  \fill[ubNavy] (current page.south west) rectangle (current page.north east);
  \fill[ubAccent] (current page.north west) rectangle ([xshift=0.10cm]current page.south west);

  \node[anchor=north west, text width=13.6cm, align=left, inner sep=0pt]
    at ($(current page.north west)+(1.00cm,-0.95cm)$) {%
      {\color{ubAccent}\bfseries\fontsize{8.2}{10}\selectfont ASESMEN MINGGU INI}\\[6pt]
      {\color{white}\bfseries\fontsize{24}{28}\selectfont #1}%
    };

  \node[anchor=north west, text width=13.6cm, align=left, inner sep=0pt]
    at ($(current page.north west)+(1.00cm,-4.55cm)$) {%
      {\color{ubNavy!0!ubHead}\rule{13.4cm}{0.4pt}}\\[7pt]
      {\color{ubAccent}\bfseries\fontsize{7.4}{9}\selectfont
       PENGAIT KE MATERI BERIKUTNYA}\\[4pt]
      {\color{ubHead}\fontsize{10.6}{13.2}\selectfont #2}%
    };

  \node[anchor=south east, inner sep=0pt]
    at ($(current page.south east)+(-0.60cm,0.32cm)$)
    {\color{ubHead}\fontsize{7.4}{9}\selectfont \insertframenumber};
\end{tikzpicture}
\end{frame}%
}



\renewcommand{\materino}{09}

\begin{document}

\deckcover{09}{9}
  {Integrasi FSM dan A* (Mini Proyek AI)}
  {2.4}
  {Python / Pygame}
  {Bab XI buku ajar}
  {Blok II \textperiodcentered\ Materi 6--9}

% --------------------------------------------------------- peta
\petapertemuan{2}
  {Materi 8 \textperiodcentered\ Pathfinding Grid}
  {Materi 9 \textperiodcentered\ Integrasi FSM dan A* (Mini Proyek AI)}
  {Materi 10 \textperiodcentered\ Pengenalan Unity: Project, Scene, dan
   GameObject}
  {Mini proyek AI kelompok: kode Python/Pygame, ringkasan desain, dan laporan
   \emph{playtest} atas tiga skenario uji, dipresentasikan 10 menit.}

% --------------------------------------------------------- capaian
\begin{frame}{Capaian Pertemuan}
\setbeamerfont{itemize/enumerate body}{size=\scriptsize}
\awalan{Sub-CPMK dan indikator dari RPS, serta tujuan Bab XI buku ajar.}
\begin{columns}[T,onlytextwidth]
  \begin{column}{0.52\textwidth}
    \begin{beamercolorbox}[wd=\linewidth,sep=5pt,rounded=true]{block title}
      {\fontsize{6.8}{8}\selectfont\color{ubHead}SUB-CPMK 2.4}\\[2pt]
      {\footnotesize Mengevaluasi perilaku NPC hasil kombinasi FSM dan A*
       terhadap \emph{skenario uji} permainan.}
    \end{beamercolorbox}
    \vspace{0.14cm}

    {\fontsize{7}{8.5}\selectfont\bfseries\color{ubAccent}INDIKATOR}\par
    \vspace{0.06cm}
    \begin{itemize}
      \item Menyusun skenario uji dan mencatat hasil pengamatan.
      \item Menilai kualitas perilaku NPC: kelancaran, daya tanggap, dan beban
            komputasi.
      \item Merumuskan perbaikan berdasarkan hasil evaluasi.
    \end{itemize}
    \vspace{0.10cm}
    \begin{kotak}
      Bab ini menutup CPMK 2 dan merupakan bab bertingkat tertinggi pada blok
      Python, yaitu ranah C5: menilai, bukan menulis kode baru.
    \end{kotak}
  \end{column}
  \begin{column}{0.45\textwidth}
    {\fontsize{7}{8.5}\selectfont\bfseries\color{ubAccent}TUJUAN BAB XI}\par
    \vspace{0.08cm}
    \begin{itemize}
      \item Menghubungkan lapisan keputusan dan lapisan navigasi dalam satu NPC.
      \item Menyimpan hasil pencarian jalur dan menentukan kapan perlu diulang.
      \item Menggerakkan NPC mengikuti jalur berupa rangkaian petak.
      \item Menilai mutu perilaku dari kelancaran, daya tanggap, dan beban
            perhitungan.
      \item Menyusun skenario uji dan mencatat hasil pengamatannya.
      \item Merumuskan perbaikan beserta cara mengukur keberhasilannya.
    \end{itemize}
  \end{column}
\end{columns}
\end{frame}

% ========================================================= bagian 1
\section{Menghubungkan Dua Lapisan}
\sectionframe{01}{Menghubungkan Dua Lapisan}
  {Niat dari mesin keadaan, rute dari pencarian jalur}

\begin{frame}{Niat dan Rute}
\awalan{Setiap lapisan menjawab satu pertanyaan yang berbeda, dan keduanya
diperlukan oleh NPC yang sama.}
\begin{itemize}
  \item \textbf{Mesin keadaan} (Bab IX) menghasilkan \textbf{niat}: jawaban atas
        ``apa yang sebaiknya saya lakukan?''. Jawabannya berupa keadaan,
        misalnya mengejar.
  \item \textbf{Pencarian jalur} (Bab X) menghasilkan \textbf{rute}: jawaban atas
        ``bagaimana cara menuju ke sana?''. Jawabannya berupa rangkaian petak.
  \item Pada Bab IX, niat mengejar dilaksanakan dengan bergerak lurus ke arah
        pemain. Pada bab ini, niat yang sama dilaksanakan dengan memanggil
        pencarian jalur terlebih dahulu.
  \item Perubahan itu hanya menambah satu langkah pada fungsi tindakan;
        mesin keadaannya tidak berubah sama sekali.
\end{itemize}
\vspace{0.1cm}
\begin{kotak}
Hasil \emph{integrasi} kedua lapisan ini bukan sekadar penjumlahan: keputusan
menentukan kapan rute dicari, sedangkan rute menentukan bagaimana keputusan itu
diwujudkan sebagai gerak.
\end{kotak}
\end{frame}

\begin{frame}{Arsitektur Integrasi}
\centering
\adjustbox{max width=0.78\linewidth,max height=0.47\textheight}{%
  % chapters/11/figures/arsitektur-integrasi.tikz
% Hubungan mesin keadaan dan pencarian jalur di dalam satu NPC.
\begin{tikzpicture}[
  kotak/.style={
    rectangle, rounded corners=2pt, draw=ubPrimary, thick, fill=ubAlt,
    align=center, text width=36mm, inner sep=4pt, minimum height=11mm,
    font=\footnotesize
  },
  simpan/.style={kotak, draw=ubGreen, fill=ubGreenBg},
  panah/.style={-{Stealth[length=2.4mm]}, thick, draw=ubAccent},
  label/.style={font=\scriptsize, align=center}
]

% Lapisan keputusan.
\node[kotak] (fsm) at (0,0)
  {\textbf{Mesin keadaan}\\[2pt]\scriptsize berpatroli, mengejar, menyerang};

% Lapisan navigasi.
\node[kotak, below=16mm of fsm] (cari)
  {\textbf{Pencarian jalur A*}\\[2pt]\scriptsize dari petak sekarang ke petak pemain};

% Simpanan jalur dan pengikut jalur.
\node[simpan, right=14mm of cari] (jalur)
  {\textbf{Simpanan jalur}\\[2pt]\scriptsize rangkaian petak beserta\\waktu hitung ulang terakhir};
\node[kotak, above=16mm of jalur] (ikuti)
  {\textbf{Pengikut jalur}\\[2pt]\scriptsize bergerak menuju petak\\berikutnya pada jalur};

\draw[panah] (fsm) -- node[label, left, xshift=-2pt]
  {niat:\\kejar} (cari);
\draw[panah] (cari) -- node[label, below, yshift=-1pt]
  {jalur} (jalur);
\draw[panah] (jalur) -- node[label, left, xshift=-2pt]
  {dipakai} (ikuti);
\draw[panah] (ikuti) -- node[label, right, xshift=2pt]
  {gerak} (fsm);

\node[label, anchor=north, text width=140mm] at (3.4,-4.4)
  {Mesin keadaan menentukan niat, sedangkan pencarian jalur menentukan rute.
   Karena pencarian jalur mahal, hasilnya disimpan dan dipakai berkali-kali;
   pencarian diulang hanya menurut selang waktu tertentu atau ketika jalurnya
   sudah habis dilalui.};
\end{tikzpicture}
}
\vspace{0.10cm}
\begin{columns}[T,onlytextwidth]
  \begin{column}{0.235\textwidth}
    {\scriptsize\textbf{Mesin keadaan} --- menentukan niat dan kapan pencarian
     dijalankan.}
  \end{column}
  \begin{column}{0.235\textwidth}
    {\scriptsize\textbf{Pencarian A*} --- mencari rute pada peta berhalangan.}
  \end{column}
  \begin{column}{0.265\textwidth}
    {\scriptsize\textbf{Simpanan jalur} --- rangkaian petak beserta waktu hitung
     ulang terakhir.}
  \end{column}
  \begin{column}{0.235\textwidth}
    {\scriptsize\textbf{Pengikut jalur} --- menggerakkan NPC menuju petak
     berikutnya.}
  \end{column}
\end{columns}
\end{frame}

\begin{frame}[fragile]{Kapan Pencarian Diulang}
\awalan{Pencarian jalur mahal, sehingga hasilnya disimpan dan dipakai
berkali-kali. Pencarian hanya diulang pada dua keadaan berikut.}
\lstinputlisting[language=Python,firstline=143,lastline=149]
  {../buku_ajar/code/python/ch11_mini_proyek.py}
\vspace{0.10cm}
\begin{itemize}
  \item \textbf{Selang hitung ulang telah habis} sejak pencarian terakhir.
  \item \textbf{Jalurnya sudah habis dilalui}, misalnya NPC telah sampai di
        tujuan sementara pemain belum berpindah.
\end{itemize}
\kodejelas{Petak pertama hasil pencarian dibuang karena itulah petak tempat NPC
berada sekarang. Kedua penghitung statistik ditambah di sini, dan keduanya
menjadi bahan penilaian pada Bagian 4.}
\end{frame}

% ========================================================= bagian 2
\section{Mengikuti Jalur}
\sectionframe{02}{Mengikuti Jalur}
  {Dari rangkaian petak menjadi gerak piksel per detik}

\begin{frame}[fragile]{Pengikut Jalur}
\lstinputlisting[language=Python,firstline=180,lastline=191]
  {../buku_ajar/code/python/ch11_mini_proyek.py}
\vspace{0.10cm}
{\scriptsize
\begin{itemize}
  \item \textbf{\texttt{while}, bukan \texttt{if}.} Dalam satu bingkai NPC dapat
        melampaui lebih dari satu petak.
  \item \textbf{Toleransi \texttt{REACH\_DISTANCE}.} Tanpa toleransi, petak tidak
        pernah dibuang dan NPC berhenti selamanya.
  \item \textbf{Jalur sebagai antrean.} Petak dibuang dari depan, bukan ditelusuri
        seluruhnya setiap bingkai.
\end{itemize}}
\end{frame}

\begin{frame}{Jalur yang Menjadi Usang}
\begin{itemize}
  \item Ketika pemain berpindah, jalur yang sedang ditempuh menjadi usang:
        rangkaian petaknya masih menuju tempat pemain berada sebelumnya.
  \item NPC terus mengikutinya sampai selang hitung ulang habis, lalu
        memperbaiki arahnya.
  \item Perilaku itu wajar, dan justru itulah pertukaran yang harus dinilai:
        memperpendek selang membuat NPC lebih tanggap tetapi lebih berat
        dihitung.
  \item Memperpanjang selang menghemat perhitungan, tetapi membuat NPC sempat
        menempuh arah yang salah.
\end{itemize}
\vspace{0.1cm}
\begin{kotak}
Pada permainan sesungguhnya, keusangan sering dikurangi dengan cara yang lebih
murah: periksa apakah jalur lama masih memadai, misalnya apakah pemain masih
dekat dengan ujung jalur. Bila ya, jalur diperpanjang, bukan dihitung ulang.
\end{kotak}
\end{frame}

\begin{frame}{Gerak yang Tidak Selalu Mulus}
\begin{itemize}
  \item NPC bergerak menuju titik tengah setiap petak, sehingga lintasannya
        berupa rangkaian garis lurus antar-petak, bukan garis halus.
  \item Pada jalur yang berbelok tajam, geraknya terlihat menyudut.
  \item Cara memperhalusnya adalah \emph{path smoothing}: singkirkan petak yang
        terletak segaris dengan petak sebelumnya dan berikutnya, sehingga hanya
        titik belok yang tersisa.
\end{itemize}
\vspace{0.12cm}
\begin{kiat}
Sebelum menambahkan penghalusan, periksa lebih dahulu apakah memang diperlukan.
Pada petak yang kecil dan NPC yang sedikit, gerak menyudut hampir tidak
terlihat; penghalusan menambah perhitungan dan menambah satu hal yang dapat
salah. Ukur dulu, baru perbaiki.
\end{kiat}
\end{frame}

% ========================================================= bagian 3
\section{Batasan Kinerja dan Hitung Ulang Berkala}
\sectionframe{03}{Batasan Kinerja dan Hitung Ulang}
  {Biaya pencarian, selang \emph{repath}, dan penyebaran beban}

\begin{frame}{Menghitung Beban Perhitungan}
\awalan{Pencarian jalur adalah perhitungan termahal pada blok Python ini.}
\begin{itemize}
  \item Biaya bertambah menurut tiga hal yang \textbf{berlipat}, bukan
        menjumlah: banyaknya NPC, seringnya pencarian diulang, dan mahalnya
        satu kali pencarian.
  \item Satu kali pencarian memeriksa sekitar 90 sampai 280 petak (Bab X).
  \item Tiga NPC yang menghitung ulang setiap bingkai pada 60 bingkai per detik:
\end{itemize}
\[
  3 \times 60 \times 90 = 16.200 \text{ petak setiap detik},
\]
yaitu sekitar enam belas ribu pemeriksaan petak setiap detik, belum termasuk
menggambar, memperbarui gerak, dan menangani tabrakan.
\begin{kotak}
Hitunglah dalam satuan \textbf{per detik}, bukan per bingkai. Beban
perhitunganlah yang menentukan apakah permainan masih berjalan pada laju
bingkai yang wajar.
\end{kotak}
\end{frame}

\begin{frame}[fragile]{Kebijakan \emph{Repath} Berkala}
\lstinputlisting[language=Python,firstline=171,lastline=178]
  {../buku_ajar/code/python/ch11_mini_proyek.py}
\vspace{0.10cm}
\begin{itemize}
  \item Hasil pencarian disimpan pada NPC, lalu dipakai berkali-kali sampai
        selang waktunya habis.
  \item Selang diatur melalui pilihan \texttt{--repath}: nilai nol berarti
        pencarian diulang setiap bingkai.
  \item Memperlambat pengulangan bukan tanpa harga: selama selang itu jalurnya
        menuju tempat pemain berada sebelumnya.
\end{itemize}
\kodejelas{Kondisi \texttt{selang\_tersisa <= 0.0 or not self.jalur} memuat kedua
pemicu sekaligus. Syarat kedua mudah terlupakan, dan akibatnya NPC berhenti
bergerak walaupun pemainnya masih terlihat.}
\end{frame}

\begin{frame}{Menyebar Waktu Perhitungan}
\begin{itemize}
  \item Bila seluruh NPC menghitung ulang pada bingkai yang sama, beban
        menumpuk pada bingkai itu lalu longgar pada bingkai berikutnya.
  \item Gejalanya laju bingkai yang naik turun, walaupun rata-ratanya tetap.
  \item Penanganannya sederhana: berikan awal yang berbeda kepada setiap NPC,
        sehingga waktu perhitungan ulangnya tidak berbarengan.
\end{itemize}
\begin{center}
{\ttfamily\small self.selang\_tersisa = self.selang * (index / jumlah\_npc)}
\end{center}
\begin{itemize}
  \item NPC pertama menghitung ulang pada awal selang, NPC kedua pada sepertiga
        selang, dan seterusnya. Beban rata-ratanya sama, tetapi sebarannya rata.
\end{itemize}
\begin{kiat}
Prinsip yang sama dipakai peladen yang membagi pekerjaan latar belakang, atau
robot yang membagi waktu perhitungan penglihatannya: pekerjaan mahal yang tidak
mendesak sebaiknya disebar, bukan ditumpuk.
\end{kiat}
\end{frame}

% ========================================================= bagian 4
\section{Menilai Mutu Perilaku}
\sectionframe{04}{Menilai Mutu Perilaku}
  {Kelancaran, daya tanggap, dan beban komputasi}

\begin{frame}{Tiga Segi Penilaian}
\awalan{Inilah inti bab ini: menilai, bukan menulis. Ketiga segi diukur, bukan
dikira-kira.}
\begin{itemize}
  \item \textbf{Kelancaran} --- seberapa lancar permainan berjalan. Ukurannya
        laju bingkai yang benar-benar tercapai; laju yang turun saat NPC
        bertambah menandakan beban telah melampaui batas.
  \item \textbf{Daya tanggap} --- seberapa cepat NPC menyesuaikan diri terhadap
        perpindahan pemain. Ukurannya jarak rata-rata NPC terhadap pemain;
        jarak lebih kecil berarti NPC lebih rapat mendekati pemain.
  \item \textbf{Beban komputasi} --- seberapa banyak pekerjaan yang diperlukan
        untuk menghasilkan perilaku itu. Ukurannya banyaknya pencarian jalur
        dan banyaknya petak yang diperiksa.
\end{itemize}
\vspace{0.1cm}
\begin{kotak}
Ketiga segi itu saling bertentangan. Menambah daya tanggap berarti menambah
beban komputasi, dan beban berlebihan mengurangi kelancaran. Titik
keseimbangannya hanya dapat ditemukan dengan mengukur.
\end{kotak}
\end{frame}

\begin{frame}{Hasil Pengukuran}
{\scriptsize
\renewcommand{\arraystretch}{1.48}
\begin{tabularx}{\textwidth}{@{}>{\raggedright\arraybackslash}p{3.3cm}
    >{\raggedleft\arraybackslash}p{2.7cm}
    >{\raggedleft\arraybackslash}p{3.0cm}
    >{\raggedleft\arraybackslash}X@{}}
\toprule
\textbf{Selang hitung ulang} & \textbf{Hitung ulang} &
\textbf{Petak diperiksa} & \textbf{Jarak rata-rata} \\
\midrule
0,00 detik (setiap bingkai) & 2.700 & 251.463 & 35 piksel \\
0,15 detik & 942 & 236.940 & 34 piksel \\
0,30 detik & 858 & 240.582 & 34 piksel \\
0,60 detik & 636 & 201.051 & 38 piksel \\
1,20 detik & 594 & 189.237 & 40 piksel \\
\bottomrule
\end{tabularx}}
\sumber{Pengukuran penulis: tiga NPC selama 900 bingkai dengan jalur pemain
yang tetap, sehingga seluruh baris dapat dibandingkan secara adil. Program
menyediakan pilihan \texttt{--evaluasi} untuk menghasilkan tabel ini.}
\vspace{0.14cm}
\begin{kotak}
Jalur pemain yang tetap adalah syarat perbandingan yang adil. Tanpa itu,
perbedaan hasil dapat berasal dari jalur pemain, bukan dari selang yang diuji.
\end{kotak}
\end{frame}

\begin{frame}{Menafsirkan Hasilnya}
\begin{enumerate}
  \item \textbf{Pencarian setiap bingkai sangat boros.} Baris pertama menjalankan
        2.700 pencarian, sekitar tiga kali lipat baris ketiga, tetapi jarak
        rata-ratanya hampir sama: 35 berbanding 34 piksel. Sepertiga pekerjaan
        sudah cukup.
  \item \textbf{Selang yang terlalu panjang mulai merugikan.} Pada 1,20 detik
        bebannya paling ringan, tetapi jarak rata-ratanya membesar menjadi 40
        piksel: NPC mulai tertinggal.
  \item \textbf{Titik keseimbangannya sekitar 0,30 detik.} Beban tinggal
        sepertiga dibandingkan pencarian setiap bingkai, sedangkan jarak
        rata-ratanya justru yang terkecil. Nilai itulah yang dipakai sebagai
        bawaan program.
\end{enumerate}
\vspace{0.12cm}
\begin{peringatan}
Kesimpulan ini berlaku untuk peta, jumlah NPC, dan jalur pemain yang diuji,
bukan untuk segala keadaan. Pada peta lebih luas, selang yang lebih panjang
mungkin lebih menguntungkan. Menyebutkan keadaan pengukuran adalah bagian dari
laporan yang jujur.
\end{peringatan}
\end{frame}

% ========================================================= bagian 5
\section{Menyusun Skenario Uji dan Merumuskan Perbaikan}
\sectionframe{05}{Skenario Uji dan Perbaikan}
  {Memicu perilaku, mengukur, lalu merumuskan perbaikan}

\begin{frame}{Syarat Skenario Uji yang Berguna}
\awalan{Pengujian biasa menuntut jawaban benar atau salah; skenario perilaku
menuntut keadaan yang memicu perilaku tertentu, lalu pengamatan atas mutunya.}
\begin{enumerate}
  \item \textbf{Dapat diulang dengan hasil yang sama.} Posisi dan urutan pemain
        tetap, misalnya jalur pemain ditentukan oleh nomor bingkai.
  \item \textbf{Memicu perilaku yang hendak dinilai.} Untuk menilai perilaku
        mengejar, pemain harus terlihat cukup panjang.
  \item \textbf{Memuat ukuran yang dapat dibandingkan.} Setiap skenario
        menghasilkan angka: jarak rata-rata, banyaknya pencarian, atau waktu.
  \item \textbf{Mewakili keadaan pemakaian yang sesungguhnya.} Ruang terbuka saja
        tidak cukup; permainan juga memuat lorong sempit dan tikungan.
\end{enumerate}
\vspace{0.1cm}
\begin{kotak}
Skenario pada ruang terbuka menjadi pembanding: di sana jalur lurus sudah
mencukupi, sehingga kelemahan pada skenario lain dapat terlihat jelas.
\end{kotak}
\end{frame}

\begin{frame}{Tiga Skenario yang Diminta}
{\scriptsize
\renewcommand{\arraystretch}{1.34}
\begin{tabularx}{\textwidth}{@{}>{\raggedright\arraybackslash}p{2.7cm}
    >{\raggedright\arraybackslash}p{5.2cm}
    >{\raggedright\arraybackslash}X@{}}
\toprule
\textbf{Skenario} & \textbf{Keadaan yang dipicu} & \textbf{Yang diamati} \\
\midrule
Ruang terbuka &
  Pemain berlari lurus menjauh, tanpa rintangan di antaranya. &
  Apakah NPC dapat menyusul, dan berapa lama waktunya. Hasilnya menjadi
  pembanding bagi kedua skenario berikutnya. \\
\addlinespace
Labirin &
  Pemain berlari mengelilingi lorong, sehingga jalan lurus selalu terhalang. &
  Apakah NPC menemukan jalan memutar, apakah jalurnya berubah-ubah tanpa sebab,
  dan berapa banyak pencarian yang dijalankan. \\
\addlinespace
Rintangan dinamis &
  Sebuah benda bergerak menutupi lorong yang sedang dilalui NPC. &
  Apakah NPC mengatasi jalur yang menjadi usang di tengah jalan, dan berapa lama
  memperbaikinya. Skenario ini mengungkap kelemahan terpenting rancangan. \\
\bottomrule
\end{tabularx}}
\sumber{Ketiga skenario ini diminta RPS dan harus ada pada laporan
\emph{playtest} mini proyek.}
\vspace{0.14cm}
\begin{kotak}
Ketiganya dipilih karena memicu kesulitan yang berbeda, dan ketiganya dapat
dibangun dengan perkakas yang telah tersedia. Ruang terbuka menjadi pembanding
bagi kedua skenario lainnya.
\end{kotak}
\end{frame}

\begin{frame}{Merumuskan Perbaikan}
{\scriptsize
\renewcommand{\arraystretch}{1.26}
\begin{tabularx}{\textwidth}{@{}>{\raggedright\arraybackslash}p{3.5cm}
    >{\raggedright\arraybackslash}p{4.4cm}
    >{\raggedright\arraybackslash}X@{}}
\toprule
\textbf{Temuan} & \textbf{Usulan perbaikan} & \textbf{Cara mengukur
keberhasilan} \\
\midrule
Pencarian setiap bingkai tidak membawa perbaikan perilaku &
  Pakai selang 0,3 detik sebagai bawaan. &
  Pencarian turun menjadi sepertiga, sedangkan jarak rata-rata tidak membesar. \\
\addlinespace
Seluruh NPC menghitung ulang pada bingkai yang sama, laju bingkai naik turun &
  Sebar waktu perhitungan dengan awal berbeda tiap NPC. &
  Simpangan laju antar-bingkai menyusut, rata-ratanya tidak berubah. \\
\addlinespace
Jalur dihitung ulang walaupun jalur lama masih memadai &
  Periksa apakah ujung jalur masih dekat dengan pemain; bila ya, perpanjang
  jalur alih-alih menghitung dari awal. &
  Pencarian turun tanpa membesarnya jarak rata-rata. \\
\addlinespace
Gerak NPC menyudut pada tikungan &
  Singkirkan petak yang segaris, sisakan titik belok saja. &
  Pengamatan atas gerak, dan tambahan waktu perhitungan yang masih diterima. \\
\bottomrule
\end{tabularx}}
\vspace{0.10cm}
{\footnotesize\color{ubMuted}\itshape Tiga perbaikan pertama dapat dibuktikan
dengan angka; perbaikan keempat menuntut pengamatan yang lebih halus. Dahulukan
yang dapat diukur.}
\end{frame}

% ========================================================= bagian 6
\section{Menutup Blok Python}
\sectionframe{06}{Menutup Blok Python}
  {Mini proyek, perkakas yang lengkap, dan yang menyusul}

\begin{frame}{Mini Proyek yang Dihasilkan}
\awalan{Mini proyek AI menuntut tiga keluaran, dan ketiganya telah tersedia pada
bab ini.}
\begin{itemize}
  \item \textbf{Kode program} --- mini permainan dengan NPC bermesin keadaan
        yang mencari jalurnya sendiri, tersusun atas \texttt{ch09\_fsm.py},
        \texttt{ch10\_astar.py}, dan \texttt{ch11\_mini\_proyek.py}.
  \item \textbf{Ringkasan desain} --- diagram keadaan beserta diagram arsitektur
        integrasi pada frame sebelumnya.
  \item \textbf{Laporan \emph{playtest}} --- catatan ketiga skenario uji beserta
        perbaikan yang dirumuskan.
  \item \textbf{Waktu pengumpulan} --- dikumpulkan dan dipresentasikan pada
        minggu kesembilan; Bab VIII sampai Bab XI juga menjadi bahan UTS pada
        minggu kesepuluh.
\end{itemize}
\vspace{0.1cm}
\begin{kotak}
Rubrik Lampiran A menimbang implementasi FSM 25\%, pencarian jalur 25\%,
penggabungan keduanya 20\%, pengujian dan laporan 20\%, serta presentasi dan
kerja sama 10\%. Bobot CPMK 2 adalah 10\% dari nilai akhir mata kuliah.
\end{kotak}
\end{frame}

\begin{frame}{Perkakas Blok Python yang Telah Lengkap}
\awalan{Seluruh perkakas Blok I dan Blok II kini tersedia, dan masing-masing
tetap berlaku setelah pekerjaan berpindah ke Unity.}
\begin{columns}[T,onlytextwidth]
  \begin{column}{0.49\textwidth}
    {\scriptsize
    \begin{itemize}
      \item Lingkungan kerja dan repositori --- Bab III.
      \item Siklus permainan dan pengendalian waktu --- Bab IV.
      \item Vektor, masukan, dan gerak --- Bab V.
      \item Gravitasi, lompatan, dan pantulan --- Bab VI.
      \item \texttt{Sprite}, \texttt{Group}, dan tabrakan --- Bab VII.
    \end{itemize}}
  \end{column}
  \begin{column}{0.49\textwidth}
    {\scriptsize
    \begin{itemize}
      \item Perilaku mengejar --- Bab VIII.
      \item Mesin keadaan --- Bab IX.
      \item Pencarian jalur A* --- Bab X.
      \item Penggabungan keduanya dan penilaian mutu --- Bab XI.
    \end{itemize}}
  \end{column}
\end{columns}
\vspace{0.16cm}
\begin{kotak}
Yang berubah pada Blok III bukan gagasannya, melainkan letak tanggung jawabnya:
siklus permainan, fisika, tabrakan, dan penyajian diserahkan kepada mesin.
\end{kotak}
\end{frame}

\begin{frame}{Apa yang Menyusul}
\begin{itemize}
  \item Capaian ini bukan perkakas yang ditinggalkan. Pada Blok III dan Blok IV
        pekerjaan berpindah ke Unity, dan seluruh gagasan tetap berlaku.
  \item Waktu Anda beralih kepada pekerjaan yang tidak dapat dikerjakan mesin:
        menentukan desain level, menyusun pengalaman bermain, dan menilai mutu
        hasilnya.
  \item Bekal mekanisme yang Anda bangun sendiri membuat peralihan berjalan
        mulus: ketika perilaku terasa janggal, Anda dapat menelusuri apakah
        komponennya keliru, mesinnya bekerja tak terduga, atau rancangan
        levelnya yang kurang tepat.
\end{itemize}
\vspace{0.1cm}
\begin{kiat}
Simpanlah catatan praktikum sejak Bab III. Catatan itulah bahan refleksi yang
paling jujur ketika Bab XVII membandingkan pendekatan Python dan Unity.
\end{kiat}
\end{frame}

% --------------------------------------------------------- ketentuan
\begin{frame}{Ketentuan Mini Proyek AI}
\awalan{Ketentuan ini diambil dari RPS dan menjadi dasar penilaian asesmen
minggu ini.}
\begin{columns}[T,onlytextwidth]
  \begin{column}{0.48\textwidth}
    {\scriptsize\bfseries\color{ubAccent}PELAKSANAAN DAN KELUARAN}\par
    \vspace{0.06cm}
    \begin{itemize}
      \item Kelompok beranggotakan 2--3 mahasiswa; dikumpulkan Minggu 9.
      \item Satu atau lebih NPC dengan FSM minimal \texttt{idle},
            \texttt{patrol}, dan \texttt{chase}.
      \item A* dipanggil saat \emph{state chase}, pada peta berhalangan.
      \item Keluaran: kode Python/Pygame, ringkasan desain (\emph{GDD}
            singkat), dan laporan \emph{playtest}.
    \end{itemize}
  \end{column}
  \begin{column}{0.48\textwidth}
    {\scriptsize\bfseries\color{ubAccent}PENGUJIAN DAN PRESENTASI}\par
    \vspace{0.06cm}
    \begin{itemize}
      \item Minimal tiga skenario uji: ruang terbuka, labirin, dan rintangan
            dinamis sederhana.
      \item Laporan \emph{playtest} 2--3 halaman berisi hasil pengamatan.
      \item Demonstrasi 10 menit, dilanjutkan tanya jawab dan umpan balik
            kelompok lain.
      \item Setiap anggota harus memahami bagiannya dan mampu menjawab
            pertanyaan penguji.
    \end{itemize}
  \end{column}
\end{columns}
\vspace{0.14cm}
\begin{kotak}
Presentasi dinilai dari kejelasan demonstrasi, kerja sama kelompok, dan
kemampuan menjawab, bukan dari panjang kode. Setengah nilai berada pada
penggabungan, pengujian, dan penyajiannya.
\end{kotak}
\end{frame}

% --------------------------------------------------------- aktivitas
\begin{frame}{Aktivitas di Kelas \textperiodcentered\ Kerja Kelompok}
\awalan{Pengalaman belajar minggu 9 pada RPS, dikerjakan berurutan di kelas lalu
diselesaikan sebelum presentasi.}
\begin{enumerate}
  \item Ketik ulang \texttt{ch11\_mini\_proyek.py} ke folder \texttt{code};
        berkas itu memakai \texttt{ch10\_astar.py}, sehingga keduanya harus
        satu folder.
  \item Jalankan \texttt{python ch11\_mini\_proyek.py --evaluasi}, catat
        tabelnya, lalu bandingkan dengan tabel pengukuran pada Bab XI.
  \item Uji ruang terbuka dan labirin: catat jarak rata-rata, banyaknya
        pencarian, dan kesan mutu perilaku NPC.
  \item Tambahkan satu benda bergerak yang menutupi lorong, lalu amati perilaku
        NPC ketika jalurnya menjadi usang.
  \item Terapkan satu perbaikan dari tabel perbaikan, ukur hasilnya dengan cara
        yang tercantum, lalu susun laporan \emph{playtest}.
  \item Siapkan demonstrasi 10 menit dan simpan seluruh pekerjaan ke repositori
        Git dengan pesan yang menerangkan isinya.
\end{enumerate}
\vspace{0.10cm}
{\footnotesize\color{ubAccent}\itshape Rincian lengkap: Mini Proyek AI pada
Bab XI buku ajar dan spesifikasi pada RPS minggu 9.}
\end{frame}

% --------------------------------------------------------- rangkuman
\begin{frame}{Rangkuman}
\begin{itemize}
  \item Mesin keadaan menghasilkan niat, pencarian jalur menghasilkan rute;
        menghubungkan keduanya hanya menambah satu langkah pada keadaan
        mengejar.
  \item Hasil pencarian disimpan dan dipakai berkali-kali; pencarian diulang
        ketika selang habis atau ketika jalurnya sudah dilalui.
  \item Pengikut jalur memakai perulangan \texttt{while} dan toleransi jarak;
        tanpa keduanya NPC meninggalkan petak atau berhenti selamanya.
  \item Selama selang hitung ulang jalurnya menjadi usang: itulah harga atas
        penghematan perhitungan.
  \item Menghitung ulang setiap bingkai sangat boros; titik keseimbangannya
        sekitar 0,30 detik, dan beban sebaiknya disebar antar-NPC.
  \item Mutu dinilai dari kelancaran, daya tanggap, dan beban komputasi; skenario
        uji yang berguna dapat diulang dan memuat ukuran yang dibandingkan.
\end{itemize}
\end{frame}

% --------------------------------------------------------- cek
\begin{frame}{Cek Pemahaman}
\begin{enumerate}
  \item Mengapa pencarian jalur tidak dijalankan setiap bingkai, dan sebutkan
        dua syarat yang memicunya.
  \item Sebuah pencarian memeriksa 200 petak. Hitunglah beban perhitungannya
        bila empat NPC menghitung ulang setiap bingkai pada 60 bingkai per
        detik.
  \item Mengapa jalur pemain pada skenario uji harus sama ketika membandingkan
        dua rancangan?
\end{enumerate}
\vspace{0.35cm}
{\footnotesize\color{ubMuted}\itshape Jawab tanpa membuka catatan, lalu
bandingkan dengan kelompok di sebelah Anda.}
\end{frame}

% --------------------------------------------------------- penutup
\penutup{Presentasi dan Laporan Mini Proyek AI}
  {Setelah UTS, Blok III dimulai dengan Unity: Materi 10 \textperiodcentered\
   Pengenalan Unity: \emph{Project}, \emph{Scene}, dan \emph{GameObject}.
   Perkakas Python yang telah lengkap tetap berlaku; yang berubah hanyalah
   letak tanggung jawabnya.}

\end{document}
